These checks show that the pipeline works and how expensive it is. They are \emph{not} tuned results: they used
500 demonstrations per task (a quarter of \code{config.DATA}), a CPU only, and a single seed.

\paragraph{Integration with the real SLM.} \code{SLMBackbone} was loaded with \code{Qwen/Qwen3-0.6B} and run forward
and backward through stage~3's optimiser on 8 steps. The prompt is 22 tokens long (with the 4 soft tokens), close to the
25 assumed in \S\ref{sec:compute}. The trainable groups are: encoder 117.7\,k, mission GRU 62.8\,k, memory gate 8.8\,k,
attention 57.9\,k, adapter 68.6\,k, LoRA 409.6\,k, action head 8.2\,k and value head 1.0\,k parameters. LoRA receives
non-zero gradients and none of the frozen SLM weights receive any.

\paragraph{Stage 1 (no SLM).}
\begin{center}
\small
\begin{tabularx}{\linewidth}{@{}Xccccc@{}}
\toprule
Run & Updates & Val CE & Val acc. & Write rate & Val success \\
\midrule
A: goto + pickup, lr $3\times10^{-4}$, batch 32, 12 epochs & 375 & 0.81 & 0.66 & 0.13 & 0.30 / 0.12 \\
B: goto only, lr $10^{-3}$, batch 8, 25 epochs & 1\,575 & 0.67 & 0.73 & 0.09 $\to$ 0.04 & 0.50 \\
\midrule
Random policy & -- & -- & -- & -- & 0.31 / 0.13 \\
\bottomrule
\end{tabularx}\\[2pt]
\footnotesize Val CE is the best epoch; val success is greedy, goto / pickup, after the last epoch (100 episodes per task for A, 200 for B).
\end{center}
Observations:
\begin{itemize}
  \item \textbf{The defaults train too slowly.} With \code{config.BC} defaults and 500 demonstrations per task, an
  epoch is only $\approx$31 updates, and after 12 epochs the policy is still at the random floor.
  \item \textbf{The model can fit the data, but overfits a small demo set.} With more updates the training
  cross-entropy falls to 0.10 while the validation cross-entropy rises after epoch~9. The final greedy success on goto
  is 0.50 (200 val episodes), well above the random 0.31 but far from the expert's 1.0. Use all 2\,000 demonstrations
  per task, select checkpoints on val, and consider DAgger.
  \item \textbf{The gate starts to collapse.} The write rate drops from 0.09 to 0.04 as training goes on. This is the
  failure mode described in \S\ref{sec:memory}: these demonstrations reward little memory, so the penalty and the
  negative bias win. Run one stage-1 variant with \code{write\_penalty = 0}.
  \item \textbf{Timing.} Stage~1 takes about 1 minute per epoch on a 22-core CPU, mostly spent in the 100-episode
  validation rollouts. Batched validation would speed this up considerably.
\end{itemize}
